% ── §15 Results: Reconstruction Assessment ──
\section{Results: Structure Reconstruction Assessment}
\label{sec:results-reconstruction}

\subsection{Question}
Can Quine--McCluskey-minimised contact circuits supply enough correct spatial
constraints to determine a protein fold?  We answer this by measuring recall
(the fraction of native contacts recovered at various prediction depths) and
long-range precision (the standard DCA evaluation metric), then comparing
against literature thresholds for structure reconstruction readiness.

\subsection{Recall Curves}
For each held-out protein and each method, we compute the number of true
contacts among the top-$K$ ranked pairs for $K \in \{L/10, L/5, L/2, L\}$,
then divide by the total native contact count:

\begin{table}[ht]
	\centering
	\caption{Mean recall over held-out proteins (100/49 split).}
	\label{tab:recall}
	\begin{tabular}{lcccc}
		\toprule
		Method           & recall@L/10 & recall@L/5 & recall@L/2 & recall@L \\
		\midrule
		Circuit-L2       & .001        & .003       & .013       & .028     \\
		Plain MI         & ---         & ---        & ---        & ---      \\
		PSICOV published & .037        & .071       & .143       & .218     \\
		\bottomrule
	\end{tabular}
\end{table}

Even at top-$L$, the circuits recover only 2.8\% of the native contact graph,
compared to DCA's 21.8\%.  The absolute number of correct constraints supplied
by the circuits is therefore far below what any contact-guided folding method
requires.

\subsection{Long-Range Precision}
The standard DCA evaluation restricts attention to long-range contacts
($|i - j| \geq 24$).  On this subset, the median precision@L/5 is:
\begin{itemize}[nosep]
	\item Circuit-L2: 0.000 (most proteins have zero long-range hits in top-L/5)
	\item Plain MI: variable, typically below 0.05
	\item PSICOV published: \textbf{0.683} (consistent with Jones et al.'s report
	      that 118/150 targets achieve $\geq 0.5$)
\end{itemize}

\subsection{Reconstruction Readiness Verdict}
We define a method as reconstruction-ready if its median long-range
precision@L/5 $\geq 0.30$ AND recall@L $\geq 0.10$ (a conservative reading of
the thresholds implied by Jones et al.\ and the CONFOLD line of work).  Under
this criterion:
\begin{center}
	\begin{tabular}{lc}
		\toprule
		Method           & Ready?       \\
		\midrule
		Circuit-L2       & No           \\
		Plain MI         & No           \\
		Perplexity       & No           \\
		Separation-only  & No           \\
		PSICOV published & \textbf{Yes} \\
		\bottomrule
	\end{tabular}
\end{center}

\subsection{Interpretation}
At current signal strength, Boolean contact circuits serve as an annotation
layer that identifies specific coevolutionary constraints; they do not supply
sufficient spatial information for de novo structure determination.  The gap
to DCA-level performance reflects the absence of phylogenetic correction and
global coupling decomposition in the circuit features---limitations that are
inherent to the consensus-residue encoding rather than artifacts of the
Quine--McCluskey minimisation itself.
